

Petri nets are a well-known model of a class of distributed systems;
for an introduction, we refer, \eg, to the monographs~\cite{reisig2013} and~\cite{BestDevillers2024}.

The \emph{reachability} problem asks whether a given target marking is reachable from a given initial marking. Here, a \emph{marking} is the standard term for a Petri net configuration, assigning a number of \emph{tokens} to each \emph{place} (local state). The reachability problem for Petri nets is well known for its computational complexity: its Ackermann-completeness has only recently been established (see~\cite{DBLP:conf/focs/Leroux21,DBLP:conf/focs/CzerwinskiO21} for the lower bound and~\cite{DBLP:conf/lics/LerouxS19} for the upper bound).

Given a Petri net and an initial marking, the \emph{boundedness} problem asks whether the reachability set is finite, and the \emph{liveness} problem asks whether all \emph{transitions} (system actions) remain live forever. Both are among the standard analysis problems. While boundedness is known to be EXPSPACE-complete~\cite{DBLP:conf/stoc/CardozaLM76,DBLP:journals/tcs/Rackoff78}, liveness is tightly related to reachability~\cite{DBLP:conf/focs/Hack74} and is thus now known to be Ackermann-complete.


There are also natural structural versions of boundedness and liveness.
The \emph{structural boundedness} problem asks, given a Petri net, whether the net is bounded for every initial marking, while the \emph{structural liveness} problem asks, given a Petri net, whether there exists an initial marking for which the net is live.

While structural boundedness is easily shown to be in PTIME (and is thus substantially easier than boundedness), structural liveness is only known to be EXPSPACE-hard and decidable~\cite{DBLP:journals/acta/JancarP19}. In fact, the decidability result can be strengthened by recent results on the home-space problem~\cite{DBLP:conf/concur/JancarL23,DBLP:journals/lmcs/JancarL24}, which readily imply an Ackermannian upper bound for structural liveness as well; nevertheless, the large complexity gap still calls for clarification.


\emph{Our contribution.}
As a step towards clarifying the complexity of \emph{structural
liveness} for general Petri nets, we establish
\emph{EXPSPACE-completeness} for the class of \emph{conservative
nets}, which preserve a weighted sum of tokens throughout their
executions. We recall that deciding whether a given net is
conservative is also in PTIME, similarly to structural boundedness
(see, \eg,~\cite{DBLP:conf/apn/MayrW14}).

A crucial notion in our proof is \emph{structural reversibility},
referred to simply as \emph{reversibility} in this paper. 
A net is reversible if there exists a sequence of transitions in which
each transition occurs at least once and whose overall effect is zero,
that is, executing the sequence does not change the marking.
Equivalently, the effect of any transition can be ``undone'' by some transition sequence.
Reversibility is also in PTIME, and it can be easily shown to be a~necessary condition for structural liveness in the case of structurally bounded nets.

Moreover, it is immediate that every conservative net is structurally bounded, and a straightforward application of Farkas’ lemma shows that every reversible structurally bounded net is conservative. Hence, our EXPSPACE-completeness result can equivalently be stated for the class of \emph{structurally bounded nets}.

As a natural first step in our future research, we plan to address several subtle issues that would allow us to extend the result to the full class of reversible nets.

The lower bound, EXPSPACE-hardness, is achieved by adapting the
construction of~\cite{DBLP:journals/acta/JancarP19}, which shows a
reduction from the EXPSPACE-complete word problem for commutative
semigroups. This problem can also be phrased as a coverability problem
for reversible Petri nets~\cite{DBLP:conf/stoc/CardozaLM76,
mayr1982complexity, DBLP:journals/jc/Mayr97}. We recall that
coverability is a weaker form of reachability: it asks whether there
exists a reachable marking that is component-wise at least as large as the target. Our adaptation demonstrates EXPSPACE-hardness of structural liveness even for nets where each transition has precisely two input and 
two output places, that is, for nets that naturally correspond to population protocols~\cite{DBLP:journals/dc/AngluinADFP06}.


Our main result is an EXPSPACE upper bound. The crucial step
establishes
that for every structurally live conservative net, there exists a live
marking with at most a doubly exponential (2-exp) number of tokens.
In fact, we establish this 2-exp bound in a more general setting: for nets
that possess live bottom strongly connected components in their
reachability graphs. Conservative nets constitute a
special case, for which such a 2-exp upper bound readily implies an
EXPSPACE upper bound for the structural liveness problem.

The aforementioned 2-exp bound is achieved by showing that for any
such net, there exists a specific quantifier-free
Presburger formula that admits a ``small'' (2-exp) solution and whose
solutions are live markings. More precisely, a solution to this
formula represents a set of at most exponentially many mutually
reachable markings, which together demonstrate that all of them are live.

A key aspect of our approach is the replacement of
reachability with a conceptually simpler \emph{virtual reachability},
which allows for negative numbers of tokens in (virtual) markings. To
express virtual reachability in reversible nets, we employ linear
systems---that  is, Boolean combinations of linear equalities,
inequalities, and
divisibility constraints---which constitute the aforementioned
special
Presburger formulas. As a technical contribution of independent
interest, we provide exponential bounds on the minimal solutions of
such linear systems, extending known results such as those
in~\cite{DBLP:conf/rta/Pottier91} (see also the survey~\cite{DBLP:journals/siglog/Haase18}).




\emph{Related research.}
This paper can be viewed as a continuation of the research line
initiated by \cite{DBLP:journals/ipl/BestE16}, which explicitly
indicated that even the decidability of structural liveness for Petri
nets remained open at that time. As already mentioned, decidability
and EXPSPACE-hardness are now established
\cite{DBLP:journals/acta/JancarP19}; furthermore, the decidability
result was recently strengthened by \cite{DBLP:conf/concur/JancarL23,
DBLP:journals/lmcs/JancarL24}, implying an Ackermannian upper bound.
Another result in this research line is the PSPACE-completeness of
structural liveness for IO-nets (Immediate Observation Petri Nets),
which were introduced in \cite{DBLP:conf/apn/EsparzaRW19} and inspired
by a subclass of population protocols in
\cite{DBLP:journals/dc/AngluinAER07}. While
\cite{DBLP:conf/apn/EsparzaRW19} does not consider structural liveness
explicitly, a~PSPACE upper bound follows directly from its results on liveness; an explicit, self-contained proof of PSPACE-completeness is provided in \cite{DBLP:journals/fuin/JancarV22}.


Numerous studies have investigated liveness for various subclasses of
Petri nets, often exploring related structural properties (see
the monographs~\cite{reisig2013, BestDevillers2024,
desel_esparza_1995}, or~\cite{DBLP:journals/fuin/HujsaD18} as 
an example of a paper explicitly studying structural liveness 
for a specific subclass). In particular, we use established results
regarding liveness for conservative nets, for which we refer to~\cite{DBLP:conf/apn/MayrW14}.




\emph{Organization of the paper.} 
Section~\ref{sec:basicdef} introduces basic notions and notation,
defines linear systems, 
and states the main results; it also establishes the equivalence of conservativeness and structural boundedness for reversible nets. Section~\ref{sec:upperbound} establishes the aforementioned 2-exp upper bound, drawing on results proven in Sections~\ref{sec:smallsolutions} and~\ref{sec:virtual}. Finally, Section~\ref{sec:hardness} addresses the EXPSPACE lower bound.

We note that throughout this paper, we adopt a convention where lemmas are considered more significant than propositions.
